\documentclass[dvipdfmx,a4paper]{article}
\usepackage{amsmath}
\usepackage{amssymb}
\usepackage{amsthm}
\usepackage{amsfonts}
\usepackage{bm}
\usepackage{geometry}
\geometry{left=25mm,right=25mm,top=25mm,bottom=25mm}

\title{G(3,1,1) Extension of Double Spacetime Torsional Mismatch: 10-Dimensional Generator Structure and Plane-Based Bracket Products}

\author{https://github.com/hypernumbernet}
\date{\today}

\theoremstyle{definition}
\newtheorem{definition}{Definition}
\newtheorem{conjecture}{Conjecture}
\theoremstyle{plain}
\newtheorem{theorem}{Theorem}
\theoremstyle{remark}
\newtheorem{remark}{Remark}

\begin{document}

\maketitle

\begin{abstract}
We embed the double spacetime structure into the projective geometric algebra \(G(3,1,1)\) by adjoining an additional null vector \(e_{4}\) to the biquaternion basis of \(\mathrm{Cl}(3,1)\). The resulting torsional mismatch is generated by a ten-dimensional bivector space consisting of three hyperbolic generators (\(iI, iJ, iK\)), three cyclic generators (\(I, J, K\)), and four null generators \(N_{\mu} := e_{4}\wedge e_{\mu}\) (\(\mu=0,1,2,3\)) constructed from \(e_{4}\) and the four usual-spacetime vector generators \(\{j, kI, kJ, kK\}\). This space is a structural candidate for the Poincar\'{e} algebra \(\mathfrak{iso}(3,1)\cong\mathfrak{so}(3,1)\ltimes\mathbb{R}^{3,1}\): the six hyperbolic-cyclic bivectors span a Lorentz copy \(\mathfrak{so}(3,1)\), while the four null bivectors form an abelian translation ideal under the geometric product.

The quadratic form \(J^{(5)}\) on torsion-plus-translation parameters combines the Cartan--Killing quadratic on the six torsional coefficients with the Minkowski metric on the translation coefficients; we use it as a scalar diagnostic on parameter space. Plane-based bracket products on three-blades and the dual-as-normal correspondence are left as an open geometric programme (Appendix~B supplies representative calculations). Because every null generator \(N_{\mu}\) is a bivector that squares to zero and is reversal-odd (\(\widetilde{N_{\mu}}=-N_{\mu}\)), the motor defined by the product \(M:=RT\) with \(R=\exp(\Omega_{\mathrm{torsion}})\) and \(T=1+\Omega_{\mathrm{trans}}\) is unitary, \(M\widetilde{M}=1\). When \([\Omega_{\mathrm{torsion}},\Omega_{\mathrm{trans}}]\neq 0\), the relation between this product and \(\exp(\Omega_{\mathrm{biv}}^{(5)})\) is left for later work.
\end{abstract}
 
\section{Introduction}
\label{sec:introduction}

The double spacetime framework describes gravity and inertia as torsional mismatch between two intrinsically paired Minkowski structures carried by each massive particle. To incorporate translational degrees of freedom in a geometrically natural manner, we lift the construction into the projective geometric algebra \(G(3,1,1)\)~\cite{Gunn2017,Lengyel2024,DeKeninckDorst2019,GunnDeKeninck2020}, whose generating vector space is five-dimensional and whose full Clifford algebra is \(2^{5}=32\)-dimensional. The use of a degenerate vector with vanishing square to encode translations as bivectors goes back to Selig's treatment of rigid-body kinematics in conformal/projective Clifford algebras~\cite{Selig2005}, and is now the standard PGA convention for both Euclidean and Minkowski signatures~\cite{Lengyel2024,GunnDeKeninck2020}. The additional null vector \(e_{4}\) anticommutes with the four \(\mathrm{Cl}(3,1)\) vector generators \(\{j, kI, kJ, kK\}\)~\cite{Hestenes2015,DoranLasenby2003} and satisfies \(e_{4}^{2}=0\). It moreover commutes with the biquaternionic pseudoscalar \(i=jkI\,kJ\,kK\)~\cite{ConwaySmith2003}, so the duality map \(X\mapsto Xi\) acts on the \(\mathrm{Cl}(3,1)\) sub-algebra in the same way as in the parent theory while leaving the new null direction intact.

The mismatch bivectors partition into three distinct classes. The three hyperbolic generators \(iI, iJ, iK\) square to \(+1\) and span the boost sector of \(\mathfrak{so}(3,1)\). The three cyclic generators \(I, J, K\) square to \(-1\) and span the rotation sector of that same Lorentz copy (DST duality interchanges the two sectors up to signs). The four null generators \(N_{\mu} := e_{4}\wedge e_{\mu}\) (\(\mu=0,1,2,3\)), constructed from \(e_{4}\) and the four Cl(3,1) vector generators, square to zero and form an abelian translation ideal. This yields a ten-dimensional bivector space matching the dimension of the Poincar\'{e} algebra \(\mathfrak{iso}(3,1)\cong\mathfrak{so}(3,1)\ltimes\mathbb{R}^{3,1}\)~\cite{Weinberg1995,Tung1985}. We treat the match as a structural candidate; a complete Lie-bracket isomorphism is beyond the scope of this paper. Translations therefore emerge directly from the null sector without auxiliary gauge fields, and the resulting extended mismatch bivector unifies the original torsional degrees of freedom with translational ones inside a single algebraic object.

A second theme is the plane geometry of \(G(3,1,1)\). Plane-based bracket products on three-blades suggest a geometric reading of the dual spacetime as a normal bundle; Appendix~B records representative calculations, while a general normal-mapping theorem remains open.

Because each null bivector squares to zero \emph{and} satisfies \(N_{\mu}N_{\nu}=0\) for every \(\mu,\nu\), the exponential of the null part is a finite polynomial of degree one. Combined with the reversal-odd character \(\widetilde{N_{\mu}}=-N_{\mu}\) of all bivectors, this guarantees the unitarity \(M\widetilde{M}=1\) of the motor defined by the product \(M:=RT\). Sandwiching with respect to the full degenerate metric of \(G(3,1,1)\), and the identification of \(M\) with \(\exp(\Omega_{\mathrm{biv}}^{(5)})\) when the factors fail to commute, are left aside here.

The paper is organised as follows. Section~2 defines the embedding of the double spacetime into \(G(3,1,1)\) and classifies the ten generators. Section~3 constructs the extended mismatch bivector, proves the polynomial form of the null exponential, and introduces the parameter-space form \(J^{(5)}\). Section~4 outlines the plane-based bracket programme. Section~5 establishes unitarity of the product motor \(M:=RT\). Section~6 develops the PGA--TEGR correspondence on a fixed chart, with the Schwarzschild diagonal ansatz and the identification of \(J^{(5)}\) with the teleparallel scalar \(T\). Conclusions are drawn in Section~7.

\section{The Ten-Dimensional Generator Structure}

\subsection{Vector basis and the adjunction of \(e_{4}\)}

The parent theory identifies the four \(\mathrm{Cl}(3,1)\) vector generators with biquaternion elements
\[
e_{0}=j,\qquad e_{1}=kI,\qquad e_{2}=kJ,\qquad e_{3}=kK,
\]
satisfying the Minkowski signature \((-,+,+,+)\):
\[
e_{0}^{2}=-1,\qquad e_{1}^{2}=e_{2}^{2}=e_{3}^{2}=+1,\qquad \{e_{\mu},e_{\nu}\}=0\ (\mu\neq\nu).
\]
To incorporate translational degrees of freedom while preserving the torsional character of the double spacetime framework, we lift this structure to \(G(3,1,1)\) by adjoining a single new vector generator \(e_{4}\) characterised by
\begin{equation}
e_{4}^{2}=0,\qquad \{e_{4},e_{\mu}\}=0\quad (\mu=0,1,2,3).
\label{eq:e4-defs}
\end{equation}
The full Clifford algebra \(G(3,1,1)\) has dimension \(2^{5}=32\); the original 16-dimensional biquaternion algebra is recovered as the sub-algebra generated by \(\{e_{0},e_{1},e_{2},e_{3}\}\).

\subsection{Compatibility with the biquaternion duality map}

The biquaternionic pseudoscalar of the parent theory is
\[
i = e_{0}e_{1}e_{2}e_{3} = j\,(kI)(kJ)(kK),
\]
and the duality map of the double spacetime acts as \(X\mapsto Xi\). Using \(\{e_{4},e_{\mu}\}=0\) repeatedly,
\[
e_{4}\,i \;=\; e_{4}\,e_{0}e_{1}e_{2}e_{3} \;=\; (-1)^{4}\,e_{0}e_{1}e_{2}e_{3}\,e_{4} \;=\; i\,e_{4},
\]
so \(e_{4}\) commutes with \(i\). Consequently, the duality map \(X\mapsto Xi\) operates on the \(\mathrm{Cl}(3,1)\) sub-algebra exactly as in the parent paper, while \(e_{4}\) is left intact as an independent null direction. This is the precise sense in which the duality is extended to \(G(3,1,1)\) ``while preserving the null character of \(e_{4}\)''.

\subsection{Classification of the ten bivector generators}

The bivector space of \(G(3,1,1)\) has dimension \(\binom{5}{2}=10\). Combining the six \(\mathrm{Cl}(3,1)\) bivectors with the four bivectors involving \(e_{4}\), the generators of the extended torsional mismatch fall into three distinct classes according to the signature of their squares:

\begin{itemize}
  \item \textbf{Hyperbolic generators} (square to \(+1\)):
  \[
  iI = e_{0}e_{1},\qquad iJ = e_{0}e_{2},\qquad iK = e_{0}e_{3}.
  \]
  These three elements span the boost sector of \(\mathfrak{so}(3,1)\).

  \item \textbf{Cyclic generators} (square to \(-1\)):
  \[
  I = e_{3}e_{2},\qquad J = e_{1}e_{3},\qquad K = e_{2}e_{1}.
  \]
  These three elements span the rotation sector of the same Lorentz copy \(\mathfrak{so}(3,1)\).

  \item \textbf{Null generators} (square to \(0\)): four bivectors built from \(e_{4}\) and the four \(\mathrm{Cl}(3,1)\) vector generators,
  \[
  N_{0} = e_{4}\wedge e_{0} = e_{4}\,j,\qquad
  N_{1} = e_{4}\wedge e_{1} = e_{4}\,(kI),
  \]
  \[
  N_{2} = e_{4}\wedge e_{2} = e_{4}\,(kJ),\qquad
  N_{3} = e_{4}\wedge e_{3} = e_{4}\,(kK).
  \]
  Using \eqref{eq:e4-defs}, we have for any \(\mu,\nu\in\{0,1,2,3\}\),
  \begin{equation}
  N_{\mu}N_{\nu} \;=\; (e_{4}e_{\mu})(e_{4}e_{\nu}) \;=\; -\,e_{4}^{2}\,e_{\mu}e_{\nu} \;=\; 0,
  \label{eq:null-strong-vanishing}
  \end{equation}
  so each \(N_{\mu}\) squares to zero (the case \(\mu=\nu\)) and \emph{any pairwise product} of null generators also vanishes. The four \(N_{\mu}\) span an abelian translation ideal of the ten-dimensional algebra.
\end{itemize}

Collectively, these classes furnish a ten-dimensional bivector space:
\[
\dim \;=\; 3 \text{ (hyperbolic)} \;+\; 3 \text{ (cyclic)} \;+\; 4 \text{ (null)} \;=\; 10,
\]
matching the dimension of the Poincar\'{e} algebra \(\mathfrak{iso}(3,1)\). The Lorentz part \(\mathfrak{so}(3,1)\) is generated by the six hyperbolic-cyclic bivectors, and the four null bivectors \(N_{\mu}\) realise the translation generators \(P_{\mu}\) as an abelian ideal under the geometric product. The construction is therefore a dimensional and structural candidate for \(\mathfrak{so}(3,1)\ltimes\mathbb{R}^{3,1}\); we do not develop a complete Lie-isomorphism theorem here.

\subsection{Null sector and translations}

The null sector directly encodes spacetime translations. Indeed, for any translation parameter \(a^{\mu}\) (\(\mu=0,1,2,3\)), the combination \(T(a):=\tfrac{1}{2}a^{\mu}N_{\mu}\) (Einstein summation implied) is a bivector built from the null sector alone. By \eqref{eq:null-strong-vanishing}, \(T(a)^{2}=0\), so the exponentiation
\[
\exp\bigl(T(a)\bigr) \;=\; 1 + \tfrac{1}{2}\,a^{\mu}N_{\mu}
\]
truncates at first order. Consequently, the translational action appears as a pure shift without introducing additional curvature or higher-order terms; in particular, light-like translations correspond to the cone \(\eta_{\mu\nu}a^{\mu}a^{\nu}=0\) inside the four-dimensional parameter space.

\subsection{Duality on the null generators}

Since the duality map \(X\mapsto Xi\) acts on the \(\mathrm{Cl}(3,1)\) sub-algebra exactly as in the parent paper, it exchanges the boost generators \(\{iI,iJ,iK\}\) with the rotation generators \(\{I,J,K\}\) up to signs. Moreover, because \(e_{4}\) commutes with \(i\),
\[
N_{\mu}\,i \;=\; (e_{4}e_{\mu})\,i \;=\; e_{4}\,(e_{\mu}\,i).
\]
Here \(e_{\mu}i\) is grade~3 in the Cl(3,1) subalgebra, so \(N_{\mu}i\) is grade~4 in \(G(3,1,1)\). Consequently duality maps each null generator out of the bivector grade: the four-dimensional null bivector ideal is not dual-closed, whereas the six-dimensional torsional sector remains closed under duality within grade~2.

In summary, the embedding into \(G(3,1,1)\) organises the extended mismatch into a ten-dimensional bivector space of Poincar\'{e} type \(\mathfrak{so}(3,1)\ltimes\mathbb{R}^{3,1}\) (as a structural candidate), in which translations emerge naturally from a four-dimensional null sector built from \(e_{4}\) and the usual-spacetime vector generators.

\section{Construction of the Extended Torsional Mismatch}

We now construct the extended torsional mismatch bivector on the ten-dimensional generator space introduced in the previous section. Let \(\Omega_{\text{biv}}^{(5)}\) denote the infinitesimal generator of the extended motor. We decompose it according to the three classes of generators:

\begin{equation}
\Omega_{\text{biv}}^{(5)}
= \underbrace{\sum_{a=1}^{3} \Bigl( \frac{\alpha_{a}}{2}\, B_{a}^{+} + \frac{\beta_{a}}{2}\, B_{a}^{-} \Bigr)}_{\Omega_{\text{torsion}}\ :\ \text{hyperbolic + cyclic (6D)}}
\;+\;
\underbrace{\sum_{\mu=0}^{3} \frac{\lambda^{\mu}}{2}\, N_{\mu}}_{\Omega_{\text{trans}}\ :\ \text{null sector (4D)}},
\label{eq:Omega-decomp}
\end{equation}

where \(B_{a}^{+}\) are the hyperbolic generators (\(iI, iJ, iK\)), \(B_{a}^{-}\) are the cyclic generators (\(I, J, K\)), and \(N_{\mu}=e_{4}\wedge e_{\mu}\) (\(\mu=0,1,2,3\)) are the four null bivectors constructed in Section~2.

The coefficients \(\alpha_{a}\) and \(\beta_{a}\) parameterise the original torsional mismatch between the usual and dual sectors, while the four coefficients \(\lambda^{\mu}\) encode the translational mismatch as a Lorentz four-vector. Light-like translations correspond to the cone \(\eta_{\mu\nu}\lambda^{\mu}\lambda^{\nu}=0\) inside the four-dimensional parameter space.

\subsection{Rigorous truncation of the null exponential}

Because the strong vanishing property \eqref{eq:null-strong-vanishing} ensures \(N_{\mu}N_{\nu}=0\) for \emph{all} pairs \((\mu,\nu)\)---not merely for \(\mu=\nu\)---we have
\begin{equation}
\left(\sum_{\mu=0}^{3}\frac{\lambda^{\mu}}{2}N_{\mu}\right)^{\!\!2}
\;=\;\frac{1}{4}\sum_{\mu,\nu=0}^{3}\lambda^{\mu}\lambda^{\nu}\,N_{\mu}N_{\nu}
\;=\;0,
\end{equation}
and therefore the exponential of the null part truncates at first order:
\begin{equation}
\exp\!\Bigl( \Omega_{\text{trans}} \Bigr)
\;=\;\exp\!\Bigl(\sum_{\mu=0}^{3}\frac{\lambda^{\mu}}{2}N_{\mu}\Bigr)
\;=\;1+\sum_{\mu=0}^{3}\frac{\lambda^{\mu}}{2}N_{\mu}.
\label{eq:exp-truncation}
\end{equation}
We emphasise that the validity of \eqref{eq:exp-truncation} requires \emph{all} mixed products \(N_{\mu}N_{\nu}\) to vanish, not only the diagonal squares \(N_{\mu}^{2}=0\); the strong vanishing \eqref{eq:null-strong-vanishing} is precisely what supplies this.

\subsection{Motor factorisation}

Because every \(N_{\mu}\) satisfies \(N_{\mu}N_{\nu}=0\), the four null generators commute. The cross-commutators \([B_{a}^{\pm},N_{\mu}]\) are generally nonzero, however, so \(\Omega_{\text{torsion}}\) and \(\Omega_{\text{trans}}\) do not commute in general. We therefore work with the motor defined by the ordered product
\begin{equation}
M \;:=\; R\,T,
\qquad
R := \exp\!\bigl(\Omega_{\text{torsion}}\bigr),\qquad
T := \exp\!\bigl(\Omega_{\text{trans}}\bigr)=1+\Omega_{\text{trans}},
\end{equation}
for which Section~5 proves \(M\widetilde{M}=1\). When the factors fail to commute, \(\exp(\Omega_{\text{biv}}^{(5)})\) differs in general from \(RT\); that identification (and any conjugated re-parameterisation \(T=\exp(\Omega'_{\text{trans}})\)) is deferred.

\subsection{Duality revisited}

The hyperbolic and cyclic components transform into each other under duality in the manner of the parent paper. As recorded in Section~2, duality sends each null generator \(N_{\mu}\) to a grade~4 element, so translational parameters must be tracked separately from the dualised torsional sector.

\subsection{The quadratic invariant on the ten-dimensional space}

It is well known that the Cartan--Killing form on the Poincar\'{e} algebra \(\mathfrak{iso}(3,1)\) is \emph{degenerate}~\cite{FultonHarris1991,Humphreys1972}: for the translation generators \(P_{\mu}\),
\[
B_{\text{Killing}}(P_{\mu},P_{\nu})\;=\;0.
\]
Consequently the Killing form alone cannot supply a quadratic invariant that detects the translational mismatch. We therefore extend the Killing form of the Lorentz sector by the Minkowski metric on the translation sector, obtaining a non-degenerate bilinear form on the ten-dimensional generator space. Explicitly we set
\begin{equation}
\boxed{\;
J^{(5)}
\;:=\;
\tfrac{1}{16}\,B_{\text{Killing}}\!\bigl(\Omega_{\text{torsion}},\Omega_{\text{torsion}}\bigr)
\;+\;
\tfrac{1}{2}\,\eta_{\mu\nu}\,\lambda^{\mu}\lambda^{\nu},
\;}
\label{eq:J5-invariant}
\end{equation}
where \(B_{\text{Killing}}\) denotes the Cartan--Killing-type quadratic form on the six-dimensional torsional (Lorentz) coefficients used in the parent paper and \(\eta_{\mu\nu}=\mathrm{diag}(-1,+1,+1,+1)\) is the Minkowski metric on the translation sector. The first term reproduces the torsional scalar \(J\) of the parent theory, while the second term provides a Lorentz-invariant quadratic measure of translational mismatch which vanishes precisely on the light cone \(\eta_{\mu\nu}\lambda^{\mu}\lambda^{\nu}=0\) where translations are null. The two contributions are decoupled because the Killing form on \(\mathfrak{iso}(3,1)\) vanishes between semisimple and translation generators.

\begin{remark}[Boundedness of \(J^{(5)}\)]
\label{rem:J5-bounded}
The translational contribution \(\tfrac{1}{2}\eta_{\mu\nu}\lambda^{\mu}\lambda^{\nu}\) is unbounded when \(\lambda^{\mu}\) is unconstrained: only the null condition \(\eta_{\mu\nu}\lambda^{\mu}\lambda^{\nu}=0\) forces it to vanish. Boundedness theorems in the discrete Diophantine companion apply to the torsion part \(J\) (equivalently the normalized scalar \(J_{\rm norm}=\frac{8}{3\pi^2}J\)) on admissible configurations, and do not extend automatically to the full form \(J^{(5)}\).
\end{remark}

In this way, the extended torsional mismatch bivector \(\Omega_{\text{biv}}^{(5)}\) unifies the original six-dimensional torsional degrees of freedom with a four-dimensional translational sector inside a single ten-dimensional algebraic object, and the natural non-degenerate quadratic invariant \eqref{eq:J5-invariant} couples to both sectors.

\section{Plane-Based Bracket Products}
\label{sec:brackets-open}

\subsection{Bracket products and their grade content}

In \(G(3,1,1)\), planes are represented by three-blades. For two three-blades \(A\) and \(B\) we define the \emph{plane-based bracket products} as the antisymmetric and symmetric parts of the geometric product,
\begin{equation}
A \mathbin{\underline{\vee}} B \;:=\; \tfrac{1}{2}(AB - BA),
\qquad
A \mathbin{\overline{\vee}} B \;:=\; \tfrac{1}{2}(AB + BA),
\label{eq:bracket-def}
\end{equation}
where the products on the right-hand side are the geometric products of the algebra.\footnote{The symbols \(\underline{\vee}\) and \(\overline{\vee}\) are used here for the bracket products as defined in \eqref{eq:bracket-def}; this should not be confused with the ``anti-inner'' and ``anti-outer'' products of standard projective geometric algebra~\cite{Lengyel2024,DorstFontijneMann2007}.}

For two homogeneous multivectors of grade three in the five-vector algebra \(G(3,1,1)\), the geometric product has grade content
\[
A_{3}B_{3} \;=\; \langle A_{3}B_{3}\rangle_{0} \;+\; \langle A_{3}B_{3}\rangle_{2} \;+\; \langle A_{3}B_{3}\rangle_{4},
\]
since no grade-6 component exists in a five-dimensional vector space. Appendix~B verifies the symmetry pattern
\[
\langle A_{3}B_{3}\rangle_{0} = \langle B_{3}A_{3}\rangle_{0},
\qquad
\langle A_{3}B_{3}\rangle_{2} = -\,\langle B_{3}A_{3}\rangle_{2},
\qquad
\langle A_{3}B_{3}\rangle_{4} = \langle B_{3}A_{3}\rangle_{4}
\]
on representative examples; a theorem for arbitrary three-blades remains open.

\subsection{Dual-as-normal correspondence}

Let \(P\) be a three-blade constructed from the usual-spacetime vector generators, and let \(D\) be an element of the dual-spacetime sector. A natural conjecture is that, when \(D\) is dual to the normal direction of \(P\),
\begin{equation}
P \mathbin{\underline{\vee}} D \;=\; \kappa \, (e_{4}\wedge n),
\label{eq:normal-correspondence}
\end{equation}
for a normal vector \(n\) and coefficient \(\kappa\). Appendix~B works out a representative case; the general statement is left open.

\subsection{Translational mismatch via the anticommutator}

Expansions of \(P\mathbin{\overline{\vee}}D\) with null contributions suggest grade-4 carriers of the form \(e_{4}\wedge n\wedge e_{\mu}\). A complete calculus connecting these brackets to \(\Omega_{\text{biv}}^{(5)}\) remains future work.

\section{Unitarity of the Motor \(M:=RT\)}

\subsection{Reversal-odd character of the bivector generators}

The reversal operation \(\,\tilde{}\,\) on \(G(3,1,1)\) acts on a grade-2 element by sending it to its negative, since \((e_{\mu}e_{\nu})^{\sim} = e_{\nu}e_{\mu} = -e_{\mu}e_{\nu}\) for \(\mu\neq\nu\). In particular, every one of the ten bivector generators introduced in Section~2 is reversal-odd:
\begin{equation}
\widetilde{B_{a}^{+}} = -B_{a}^{+},\qquad
\widetilde{B_{a}^{-}} = -B_{a}^{-},\qquad
\widetilde{N_{\mu}}   = -N_{\mu}.
\label{eq:rev-odd}
\end{equation}
We stress that the reversal-odd property of \(N_{\mu}\) crucially relies on \(N_{\mu}\) being a \emph{bivector} \(e_{4}\wedge e_{\mu}\). Had we attempted to take a grade-1 null vector as the translation generator, it would have been reversal-even and the unitarity argument below would have failed.\footnote{This is the geometric reason for adopting null \emph{bivectors} rather than null \emph{vectors} as translation generators in projective geometric algebra. The same convention is used in the standard PGA literature for both Euclidean and Minkowski signatures~\cite{Gunn2017,Lengyel2024,GunnDeKeninck2020}.}

\subsection{Explicit unitarity calculation for \(M:=RT\)}

Consider the product motor \(M:=R\,T\). For the torsional rotor, the reversal-odd property \eqref{eq:rev-odd} gives the standard unitarity
\[
R \, \widetilde{R} \;=\; \exp(\Omega_{\text{torsion}})\,\exp(\widetilde{\Omega_{\text{torsion}}})
\;=\; \exp(\Omega_{\text{torsion}})\,\exp(-\Omega_{\text{torsion}}) \;=\; 1,
\]
which is the parent theory's unitarity \(R\widetilde{R}=1\) restricted to the six-dimensional torsional sector. For the null translator we use \eqref{eq:exp-truncation} together with \eqref{eq:rev-odd}:
\begin{align}
T \, \widetilde{T}
&= \left(1+\sum_{\mu}\frac{\lambda^{\mu}}{2}N_{\mu}\right)\!\left(1+\sum_{\nu}\frac{\lambda^{\nu}}{2}\widetilde{N_{\nu}}\right) \notag\\
&= \left(1+\sum_{\mu}\frac{\lambda^{\mu}}{2}N_{\mu}\right)\!\left(1-\sum_{\nu}\frac{\lambda^{\nu}}{2}N_{\nu}\right) \notag\\
&= 1 \;-\; \underbrace{\left(\sum_{\mu}\frac{\lambda^{\mu}}{2}N_{\mu}\right)^{\!\!2}}_{=\,0\ \text{by }\eqref{eq:null-strong-vanishing}}
\;=\; 1.
\label{eq:null-unitarity}
\end{align}
Because reversal is an antihomomorphism (\(\widetilde{R\,T}=\widetilde{T}\,\widetilde{R}\)),
\begin{equation}
M\,\widetilde{M}
\;=\; R\,T\,\widetilde{T}\,\widetilde{R}
\;=\; R\,\widetilde{R}
\;=\; 1.
\end{equation}
This establishes unitarity for the ordered product \(M=RT\). When \(\Omega_{\text{torsion}}\) and \(\Omega_{\text{trans}}\) fail to commute, \(M\) need not coincide with \(\exp(\Omega_{\text{biv}}^{(5)})\); sandwiching with respect to the full degenerate metric of \(G(3,1,1)\) is likewise left aside.

\subsection{Unified algebraic description}

Within this product-motor framework, light-like translations (\(\eta_{\mu\nu}\lambda^{\mu}\lambda^{\nu}=0\)) and torsional mismatch sit in the same ten-dimensional parameter space measured by \(J^{(5)}\). Because duality does not close the null bivector ideal (Section~2), translational parameters are tracked separately from dualised torsional data.

\section{PGA--TEGR Correspondence}
\label{sec:pga-tegr}
\label{sec:schwarzschild-pga}

\subsection{Scope}

The parent double-spacetime theory aims to realise teleparallel gravity: geometry is carried by torsion rather than by a non-vanishing curvature tensor. In the projective extension \(G(3,1,1)\) we develop that programme on a fixed coordinate chart. Throughout we work with the Poincar\'{e} candidate
\(\mathfrak{so}(3,1)\ltimes\mathbb{R}^{3,1}\) (and not with a twelve-dimensional label \(\mathfrak{so}(3,1)\oplus\mathfrak{so}(3,1)\)). The section constructs motor-induced tetrads, records Weitzenb\"{o}ck torsion, specialises to the exterior Schwarzschild diagonal ansatz, and states the status of the identification between the parameter-space forms \(J\), \(J^{(5)}\) and the teleparallel scalar \(T\). Full variational equivalence between TEGR and the Einstein--Hilbert action on a general manifold is taken from the TEGR literature~\cite{maluf2013} and is not re-derived here.

\begin{center}
\footnotesize
\setlength{\tabcolsep}{4pt}
\begin{tabular}{@{}p{0.42\textwidth}p{0.26\textwidth}l@{}}
\hline
Claim & Status & Formalisation \\
\hline
Motor \(M:=RT\) unitarity & proved (Sec.~5) & Lean \texttt{Motor} \\
Induced metric from diagonal Schwarzschild tetrad & proved (chart) & Lean \texttt{Gravity.Schwarzschild} \\
Weitzenb\"{o}ck \(T^{a}{}_{\mu\nu}\) and scalar \(T\) on the exterior chart & proved (chart) & Lean \texttt{Gravity.Weitzenbock} \\
Radial-boost scale factors match Schwarzschild redshift & proved & Lean \texttt{Gravity.Sandwich} \\
\(J=\tfrac12 T+\nabla\cdot B\) for the static diagonal gauge & proved (chart slice) & Lean \texttt{Gravity.Identification} \\
General \(J^{(5)}\leftrightarrow T\) for arbitrary motors & conjecture & below \\
TEGR \(\Leftrightarrow\) Einstein--Hilbert (variational) & external input & \cite{maluf2013} \\
\hline
\end{tabular}
\end{center}

\subsection{Motor-induced tetrads}

\begin{definition}[Reference coframe and sandwich]
Fix the Cl(3,1) vector generators \(\{\mathring{e}^{a}\}_{a=0}^{3}\) inside \(G(3,1,1)\). For a unitary motor \(M\) (Section~5) define the sandwiched vectors
\begin{equation}
e^{a} \;:=\; \bigl\langle M\,\mathring{e}^{a}\,\widetilde{M}\bigr\rangle_{1}.
\label{eq:motor-tetrad}
\end{equation}
On a coordinate chart \((x^{\mu})\) we write component fields \(e^{a}{}_{\mu}(x)\) and the induced metric
\begin{equation}
g_{\mu\nu} \;=\; \eta_{ab}\, e^{a}{}_{\mu}\, e^{b}{}_{\nu},
\label{eq:induced-metric}
\end{equation}
with \(\eta=\mathrm{diag}(-1,+1,+1,+1)\).
\end{definition}

When \(M=R=\exp(\Omega_{\mathrm{torsion}})\) is a pure Lorentz rotor, the sandwich realises a local Lorentz action on the reference vectors. Null translations \(T=1+\Omega_{\mathrm{trans}}\) enter the full motor \(M=RT\) but do not by themselves supply the classical Schwarzschild diagonal gauge; they remain available for translational mismatch diagnostics measured by \(J^{(5)}\).

\subsection{Weitzenb\"{o}ck torsion on a fixed chart}

\begin{definition}[Chart Weitzenb\"{o}ck torsion]
On a fixed chart, the Weitzenb\"{o}ck torsion of a coframe \(e^{a}{}_{\mu}\) is
\begin{equation}
T^{a}{}_{\mu\nu} \;=\; \partial_{\mu} e^{a}{}_{\nu} - \partial_{\nu} e^{a}{}_{\mu}.
\label{eq:weitzenbock-torsion}
\end{equation}
With inverse tetrad \(e_{a}{}^{\mu}\) and spacetime indices lowered by \(g_{\mu\nu}\), the teleparallel torsion scalar is
\begin{equation}
T \;=\; \tfrac14 T_{\rho\mu\nu}T^{\rho\mu\nu} + \tfrac12 T_{\rho\mu\nu}T^{\nu\mu\rho} - T_{\rho}T^{\rho},
\label{eq:teleparallel-T}
\end{equation}
where \(T_{\rho}=T^{\sigma}{}_{\rho\sigma}\).
\end{definition}

No curvature tensor of a metric-compatible Levi-Civita connection is required at this stage: the teleparallel description stores geometry in \(T^{a}{}_{\mu\nu}\).

\subsection{Schwarzschild ansatz}

\begin{definition}[Exterior Schwarzschild diagonal tetrad]
Let \(r_{s}=2GM/c^{2}>0\) and work in the exterior region \(r>r_{s}\). The diagonal coframe (units \(c=1\)) is
\begin{align}
e^{0} &= \sqrt{1-\frac{r_{s}}{r}}\,dt,\nonumber\\
e^{1} &= \bigl(1-\frac{r_{s}}{r}\bigr)^{-1/2}\,dr,\nonumber\\
e^{2} &= r\,d\theta,\nonumber\\
e^{3} &= r\sin\theta\,d\phi.
\label{eq:schwarzschild-tetrad}
\end{align}
\end{definition}

\begin{definition}[Radial boost rapidity]
Set \(A(r):=1-r_{s}/r\) and
\begin{equation}
\varphi(r) \;:=\; -\log\sqrt{A(r)} \;=\; -\tfrac12\log A(r)
\qquad(r>r_{s}).
\label{eq:radial-rapidity}
\end{equation}
Then \(e^{-\varphi}=\sqrt{A}\) and \(e^{\varphi}=A^{-1/2}\), so the \((t,r)\) diagonal coefficients of~\eqref{eq:schwarzschild-tetrad} are the boost scale factors of a pure hyperbolic rotor generated by \(B^{+}_{1}=e_{0}\wedge e_{1}\). The angular legs \(e^{2},e^{3}\) are taken from the spherical reference coframe. The corresponding torsional parameter is the pure boost \(\Omega_{\mathrm{torsion}}=(\varphi/2)B^{+}_{1}\), with vanishing null parameters \(\lambda^{\mu}=0\) in the diagonal gauge.
\end{definition}

\subsection{Identification of \(J\) and \(J^{(5)}\) with \(T\)}

\begin{theorem}[Chart-level statements]
On the exterior chart with~\eqref{eq:schwarzschild-tetrad}:
\begin{enumerate}
\item the induced metric~\eqref{eq:induced-metric} is the Schwarzschild metric;
\item the Weitzenb\"{o}ck torsion and the scalar \(T\) admit closed-form expressions;
\item for the radial rapidity~\eqref{eq:radial-rapidity}, the parameter-space invariant satisfies
\(J(\varphi)=\tfrac12\varphi(r)^{2}\), and the field identity
\(T=\tfrac12 T_{\mathrm{bulk}}+ \nabla_{\mu}B^{\mu}\)
holds with an explicitly constructed radial current \(B^{\mu}\) (static, diagonally gauged slice).
\end{enumerate}
\end{theorem}

\begin{conjecture}[General motor identification]
For a general smooth motor field \(M(x)=R(x)T(x)\) on a chart, the extended diagnostic \(J^{(5)}\) built from \((\Omega_{\mathrm{torsion}},\lambda)\) coincides with the teleparallel scalar~\eqref{eq:teleparallel-T} up to a total divergence:
\begin{equation}
J^{(5)} \;=\; \tfrac12\, T \;+\; \nabla_{\mu}B^{\mu}.
\label{eq:J5-T-conjecture}
\end{equation}
\end{conjecture}

The conjecture reduces to the theorem above on the static diagonal Schwarzschild slice with \(\lambda^{\mu}=0\). A proof for arbitrary motors would require a complete dictionary between \(\Omega_{\mathrm{biv}}^{(5)}\) and Weitzenb\"{o}ck contorsion beyond the chart calculus used here.

\subsection{Relation to Einstein--Hilbert}

The integral identity \(\int T\,e\,d^{4}x=-\int R\sqrt{-g}\,d^{4}x+(\text{boundary})\) and the ensuing variational equivalence with Einstein's equations are standard in TEGR~\cite{maluf2013}. We use them as external input and do not re-derive the general spacetime calculus here. Sandwiching with respect to the full degenerate metric of \(G(3,1,1)\), and the identification of \(M\) with \(\exp(\Omega_{\mathrm{biv}}^{(5)})\) when \([\Omega_{\mathrm{torsion}},\Omega_{\mathrm{trans}}]\neq0\), remain open as in Section~5.

\section{Conclusions}

We have embedded the double spacetime structure into the projective geometric algebra \(G(3,1,1)\) by adjoining a new null vector \(e_{4}\) to the biquaternion vector basis \(\{j, kI, kJ, kK\}\) of \(\mathrm{Cl}(3,1)\). The generators of the extended torsional mismatch organise themselves into a ten-dimensional bivector space consisting of three hyperbolic generators \(\{iI,iJ,iK\}\), three cyclic generators \(\{I,J,K\}\), and four null bivectors \(N_{\mu}=e_{4}\wedge e_{\mu}\) (\(\mu=0,1,2,3\)). This space is a structural candidate for \(\mathfrak{iso}(3,1)\cong\mathfrak{so}(3,1)\ltimes\mathbb{R}^{3,1}\).

The strong vanishing property \(N_{\mu}N_{\nu}=0\) for all \(\mu,\nu\) ensures that the exponential of the null part is a finite polynomial of degree one, and the reversal-odd character \(\widetilde{N_{\mu}}=-N_{\mu}\) delivers the unitarity \(M\widetilde{M}=1\) of the product motor \(M:=RT\). When the factors fail to commute, the relation to \(\exp(\Omega_{\mathrm{biv}}^{(5)})\) and sandwiching of the full degenerate metric remain open.

The parameter-space form
\[
J^{(5)} \;=\; \tfrac{1}{16}B_{\text{Killing}}\!\bigl(\Omega_{\text{torsion}},\Omega_{\text{torsion}}\bigr) \;+\; \tfrac{1}{2}\eta_{\mu\nu}\lambda^{\mu}\lambda^{\nu}
\]
serves as a scalar diagnostic on torsion-plus-translation coefficients. Duality closes the six-dimensional torsional bivector sector, while each null generator is sent to grade~4. Plane-based brackets and dual-as-normal geometry form an open programme (Appendix~B). Section~6 develops the PGA--TEGR correspondence on a fixed chart: motor-induced tetrads, Weitzenb\"{o}ck torsion, the Schwarzschild diagonal ansatz, and a chart-level link between \(J\) and \(T\), with the general \(J^{(5)}\leftrightarrow T\) identification left as a conjecture and the TEGR/Einstein--Hilbert variational equivalence taken from the literature.

\appendix

\section{Explicit Correspondence of the Ten-Dimensional Generators}

The extended torsional mismatch is generated by a ten-dimensional bivector space that decomposes into three classes. Below we give the explicit algebraic representatives in the \(G(3,1,1)\) setting, using the biquaternion vector basis \(\{e_{0}=j,\ e_{1}=kI,\ e_{2}=kJ,\ e_{3}=kK\}\) of the parent theory together with the adjoined null vector \(e_{4}\) (\(e_{4}^{2}=0\), \(\{e_{4},e_{\mu}\}=0\)).

\subsection{Hyperbolic Generators (square to \(+1\))}

These three generators are associated with boost-like transformations:

\begin{align*}
B_{1}^{+} &= iI = e_{0}\wedge e_{1}, \\
B_{2}^{+} &= iJ = e_{0}\wedge e_{2}, \\
B_{3}^{+} &= iK = e_{0}\wedge e_{3}.
\end{align*}

Each satisfies \((B_{a}^{+})^{2} = +1\).

\subsection{Cyclic Generators (square to \(-1\))}

These three generators are associated with rotation-like transformations in the dual spacetime:

\begin{align*}
B_{1}^{-} &= I = e_{3}\wedge e_{2}, \\
B_{2}^{-} &= J = e_{1}\wedge e_{3}, \\
B_{3}^{-} &= K = e_{2}\wedge e_{1}.
\end{align*}

Each satisfies \((B_{a}^{-})^{2} = -1\).

\subsection{Null Generators (square to \(0\))}

The null generators are the four bivectors built from the adjoined null vector \(e_{4}\) and the four \(\mathrm{Cl}(3,1)\) vector generators:
\begin{align*}
N_{0} &= e_{4}\wedge e_{0} = e_{4}\,j,         &\quad& \text{(time-translation generator)}\\
N_{1} &= e_{4}\wedge e_{1} = e_{4}\,(kI),      &\quad& \text{(}x\text{-translation generator)}\\
N_{2} &= e_{4}\wedge e_{2} = e_{4}\,(kJ),      &\quad& \text{(}y\text{-translation generator)}\\
N_{3} &= e_{4}\wedge e_{3} = e_{4}\,(kK),      &\quad& \text{(}z\text{-translation generator)}
\end{align*}
Using \(e_{4}^{2}=0\) and \(\{e_{4},e_{\mu}\}=0\), one finds directly that, for \emph{any} \(\mu,\nu\in\{0,1,2,3\}\),
\[
N_{\mu}\,N_{\nu} \;=\; (e_{4}e_{\mu})(e_{4}e_{\nu}) \;=\; -\,e_{4}^{2}\,e_{\mu}e_{\nu} \;=\; 0.
\]
Thus each \(N_{\mu}\) squares to zero (the case \(\mu=\nu\)) and any pairwise product also vanishes; the four \(N_{\mu}\) form an abelian translation ideal that realises the translation generators \(P_{\mu}\) of the Poincar\'{e} algebra \(\mathfrak{iso}(3,1)\). Their reversal is \(\widetilde{N_{\mu}}=-N_{\mu}\), which is essential for the unitarity of the motor (Section~5).

\paragraph{Relation to light-like vectors of the usual spacetime.}
The connection with the light-like vectors of the parent theory is direct. A general light-like vector in the usual spacetime is
\[
L \;=\; ct\,j + x\,kI + y\,kJ + z\,kK \;=\; ct\,e_{0} + x\,e_{1} + y\,e_{2} + z\,e_{3},
\]
with the Minkowski-signature null condition
\[
L^{2} \;=\; -(ct)^{2} + x^{2} + y^{2} + z^{2} \;=\; 0,
\]
in the signature \((-,+,+,+)\) of the parent paper. A translation \emph{along} this direction \(L\) is generated by the null bivector
\[
T_{L} \;:=\; e_{4}\wedge L \;=\; ct\,N_{0} + x\,N_{1} + y\,N_{2} + z\,N_{3},
\]
which is a linear combination of the four null bivectors and which itself squares to zero by the strong vanishing of the \(N_{\mu}N_{\nu}\). The light-likeness of the translation direction is thus expressed as the cone condition
\[
\eta_{\mu\nu}\,\lambda^{\mu}\lambda^{\nu} \;=\; -(ct)^{2}+x^{2}+y^{2}+z^{2} \;=\; 0
\]
on the four translation parameters \(\lambda^{\mu}=(ct,x,y,z)\), \emph{not} as an algebraic constraint on the generators themselves. The four \(N_{\mu}\) are linearly independent and the translation sector remains genuinely four-dimensional.

\subsection{Summary}

The complete set of ten generators and their algebraic properties is summarised in Table~\ref{tab:generators}.

\begin{table}[htbp]
\centering
\caption{Ten-dimensional generator structure}
\label{tab:generators}
\begin{tabular}{llcl}
\hline
Class          & Generators                                                            & Count & Square \\
\hline
Hyperbolic     & \(B_{1}^{+}=iI,\ B_{2}^{+}=iJ,\ B_{3}^{+}=iK\)                        & 3     & \(+1\) \\
Cyclic         & \(B_{1}^{-}=I,\ \ B_{2}^{-}=J,\ \ B_{3}^{-}=K\)                       & 3     & \(-1\) \\
Null           & \(N_{0}=e_{4}j,\ N_{1}=e_{4}(kI),\ N_{2}=e_{4}(kJ),\ N_{3}=e_{4}(kK)\) & 4     & \(\phantom{+}0\) \\
\hline
\end{tabular}
\end{table}

This ten-dimensional basis forms the generator space of the extended torsional mismatch bivector \(\Omega_{\text{biv}}^{(5)}\). Duality \(X\mapsto Xi\) closes the six hyperbolic-cyclic generators within the bivector grade (and \(e_{4}\) commutes with the biquaternionic pseudoscalar \(i\)), while each null generator \(N_{\mu}\) is sent to grade~4. Translations are generated intrinsically by the four-dimensional null sector built from \(e_{4}\) and the usual-spacetime vector generators, and light-like translations correspond to the cone \(\eta_{\mu\nu}\lambda^{\mu}\lambda^{\nu}=0\) in the four-dimensional translation parameter space.

\section{Details of Plane-Based Bracket Products and Dual Spacetime Correspondence}

In this appendix we provide explicit calculations that supplement the discussion in Section~4 of plane-based bracket products and the normal property of the dual spacetime.

\subsection{Definition of the bracket products and their grade content}

For two three-blades \(A\) and \(B\) in \(G(3,1,1)\), the plane-based bracket products are defined by
\[
A \mathbin{\underline{\vee}} B \;=\; \tfrac12 (AB - BA), \qquad
A \mathbin{\overline{\vee}} B \;=\; \tfrac12 (AB + BA).
\]
As explained in Section~4 (footnote), these symbols denote the bracket products in the sense of \eqref{eq:bracket-def}, and do not coincide with the ``anti-products'' of standard PGA defined via the anti-scalar unit.

In the five-vector algebra \(G(3,1,1)\), the geometric product of two three-blades has grade content
\[
A_{3}B_{3} \;=\; \langle A_{3}B_{3}\rangle_{0} \;+\; \langle A_{3}B_{3}\rangle_{2} \;+\; \langle A_{3}B_{3}\rangle_{4}.
\]
Under swap \(A\leftrightarrow B\), each grade component is expected to have definite symmetry:
\begin{equation}
\langle A_{3}B_{3}\rangle_{0} = +\langle B_{3}A_{3}\rangle_{0},
\quad
\langle A_{3}B_{3}\rangle_{2} = -\langle B_{3}A_{3}\rangle_{2},
\quad
\langle A_{3}B_{3}\rangle_{4} = +\langle B_{3}A_{3}\rangle_{4}.
\label{eq:grade-symmetry}
\end{equation}

\paragraph{Verification via an explicit example.}
Take \(A=e_{1}e_{2}e_{3}\) and \(B=e_{1}e_{2}e_{4}\) in \(G(3,1,1)\), with \(e_{1}^{2}=e_{2}^{2}=e_{3}^{2}=+1\), \(e_{4}^{2}=0\), and all generators anticommuting. A direct computation, anticommuting \(e_{3}\) past \(e_{1}e_{2}\) and using \(e_{1}e_{2}e_{1}e_{2}=-1\), yields
\[
AB \;=\; e_{1}e_{2}e_{3}\,e_{1}e_{2}e_{4} \;=\; -\,e_{3}e_{4},
\]
which is pure grade-2. Similarly,
\[
BA \;=\; e_{1}e_{2}e_{4}\,e_{1}e_{2}e_{3} \;=\; -\,e_{4}e_{3} \;=\; +\,e_{3}e_{4}.
\]
Therefore
\[
A\mathbin{\underline{\vee}}B \;=\; \tfrac{1}{2}(AB-BA) \;=\; -\,e_{3}e_{4} \;\in\; \text{grade-2},
\quad
A\mathbin{\overline{\vee}}B \;=\; \tfrac{1}{2}(AB+BA) \;=\; 0,
\]
in agreement with \eqref{eq:grade-symmetry}.

\paragraph{Verification of a non-zero grade-4 case.}
Take \(A=e_{1}e_{2}e_{3}\) and \(B=A=e_{1}e_{2}e_{3}\). Then \(AB=A^{2}=(e_{1}e_{2}e_{3})^{2}=-1\), which is grade-0; \(BA=A^{2}=-1\); so \(A\mathbin{\overline{\vee}}A=-1\) (scalar) and \(A\mathbin{\underline{\vee}}A=0\). For a grade-4 example, take \(A=e_{0}e_{1}e_{2}\) and \(B=e_{0}e_{3}e_{4}\). One has \(AB=(e_{0}e_{1}e_{2})(e_{0}e_{3}e_{4})\), whose grade-4 component is a non-zero \(e_{1}e_{2}e_{3}e_{4}\) (or related four-blade depending on the order of generators). These computations support the grade pattern on the examples considered; a general theorem for arbitrary three-blades remains open.

\subsection{Correspondence with dual spacetime normals}

Let \(P\) be a three-blade constructed from the usual-spacetime vector generators and let \(D\) be an element of the dual-spacetime sector. On representative configurations one finds
\begin{equation}
P \mathbin{\underline{\vee}} D \;=\; \kappa\,(e_{4}\wedge n),
\label{eq:appB-normal}
\end{equation}
where \(n\) is the normal vector to the plane represented by \(P\) (in the usual-spacetime sense), \(e_{4}\wedge n\) is a null bivector in the four-dimensional null sector spanned by \(\{N_{0},N_{1},N_{2},N_{3}\}\), and \(\kappa\) is a real coefficient determined by the normalisations of \(P\) and \(D\). Both sides are grade-2, consistent with \eqref{eq:grade-symmetry}.

\paragraph{Worked representative.}
Take \(P = e_{4}\wedge e_{1}\wedge e_{2}\) (a three-blade containing one \(e_{4}\)-factor and a 2-plane spanned by \(e_{1},e_{2}\)), and \(D=jI\), an element of the dual-spacetime sector of the parent theory. Using \(jI=-e_{0}e_{3}e_{2}\) (from the parent theory's grade-3 identification of the dual basis) and the anticommutation \(\{e_{4},e_{\mu}\}=0\), we obtain
\[
PD \;=\; (e_{4}e_{1}e_{2})\,(-e_{0}e_{3}e_{2}) \;=\; -\,e_{4}e_{1}e_{2}e_{0}e_{3}e_{2}.
\]
Anticommuting \(e_{2}\) past itself eliminates a factor (via \(e_{2}^{2}=+1\)) and after reordering one finds that \(PD-DP\) reduces to a non-zero multiple of \(e_{4}\wedge e_{3}\), i.e.\ a non-zero element of the null sector proportional to \(N_{3}=e_{4}\wedge e_{3}\). The direction \(e_{3}\) is precisely the unit normal to the 2-plane \(e_{1}\wedge e_{2}\) (in the spatial part), so \(\kappa\cdot(e_{4}\wedge n) = \kappa\,N_{3}\) with \(n=e_{3}\). This confirms \eqref{eq:appB-normal} for this representative.

Crucially, the right-hand side of \eqref{eq:appB-normal} is a bivector of the form \(e_{4}\wedge(\text{vector})\), i.e.\ an element of the null translation ideal. The dual spacetime element \(D\) therefore acts as the carrier of the normal direction, and the bracket \(P\mathbin{\underline{\vee}}D\) produces a generator of translations along that normal.

\subsection{Appearance of translational mismatch}

When the three-blade \(P\) is enlarged to include a contribution from the null sector,
\[
P \;\to\; P \;+\; \sum_{\mu=0}^{3}\lambda^{\mu}\,P^{(\mu)}_{\text{null}},
\qquad
P^{(\mu)}_{\text{null}} \;:=\; N_{\mu}\wedge n,
\]
the anticommutator bracket \(P\mathbin{\overline{\vee}}D\) acquires an additional term linear in the translational parameters:
\[
P \mathbin{\overline{\vee}} D
\;=\; \bigl(P\mathbin{\overline{\vee}}D\bigr)_{\text{torsional}}
\;+\; \sum_{\mu=0}^{3} c_{\mu}\,\lambda^{\mu}\,(e_{4}\wedge n\wedge e_{\mu}),
\]
where the grade-4 four-blade \(e_{4}\wedge n\wedge e_{\mu}\) is the geometric carrier of the translational mismatch. The coefficients \(c_{\mu}\) are determined by the relative orientation between the plane and the dual normal. The translational term vanishes identically when all \(\lambda^{\mu}=0\), recovering the purely torsional case. In particular, contractions of the form \(N_{\mu}\cdot e_{4}\) (which would vanish since \(e_{4}^{2}=0\)) do not appear here; the correct geometric object is the grade-4 four-blade \(e_{4}\wedge n\wedge e_{\mu}\).

\subsection{Consistency under duality}

The duality map \(X\mapsto Xi\) acts on the \(\mathrm{Cl}(3,1)\) sub-algebra in the standard way and, because \(e_{4}i=ie_{4}\), commutes through the \(e_{4}\)-factor. Dualisation of a null bivector \(N_{\mu}\) yields a grade-4 element (Section~2), so a fully duality-covariant bracket calculus on the null sector remains open.

These calculations illustrate geometric readings of torsional and translational mismatch and leave the general programme of Section~4 for later work.

\begin{thebibliography}{99}

\bibitem{Gunn2017}
C.~Gunn,
\textit{Geometric algebras for Euclidean geometry},
Advances in Applied Clifford Algebras \textbf{27} (1), 185--208 (2017).

\bibitem{Lengyel2024}
E.~Lengyel,
\textit{Projective Geometric Algebra Illuminated},
Terathon Software LLC (2024).

\bibitem{DeKeninckDorst2019}
S.~De Keninck and L.~Dorst,
\textit{Projective geometric algebra: A new framework for doing Euclidean geometry},
arXiv:1901.05873 (2019).

\bibitem{GunnDeKeninck2020}
C.~Gunn and S.~De Keninck,
\textit{Geometric algebra and computer graphics},
in: SIGGRAPH 2019 Courses, ACM, New York (2019).
See also: \texttt{https://bivector.net/}.

\bibitem{DorstFontijneMann2007}
L.~Dorst, D.~Fontijne, and S.~Mann,
\textit{Geometric Algebra for Computer Science: An Object-Oriented Approach to Geometry},
Morgan Kaufmann, San Francisco (2007).

\bibitem{DoranLasenby2003}
C.~Doran and A.~Lasenby,
\textit{Geometric Algebra for Physicists},
Cambridge University Press, Cambridge (2003).

\bibitem{Hestenes2015}
D.~Hestenes,
\textit{Space-Time Algebra}, 2nd ed.,
Birkh\"{a}user, Basel (2015).

\bibitem{HestenesSobczyk1984}
D.~Hestenes and G.~Sobczyk,
\textit{Clifford Algebra to Geometric Calculus: A Unified Language for Mathematics and Physics},
Reidel, Dordrecht (1984).

\bibitem{Lounesto2001}
P.~Lounesto,
\textit{Clifford Algebras and Spinors}, 2nd ed.,
London Mathematical Society Lecture Note Series \textbf{286},
Cambridge University Press, Cambridge (2001).

\bibitem{ConwaySmith2003}
J.~H.~Conway and D.~A.~Smith,
\textit{On Quaternions and Octonions: Their Geometry, Arithmetic, and Symmetry},
A K Peters, Natick, MA (2003).

\bibitem{Selig2005}
J.~M.~Selig,
\textit{Geometric Fundamentals of Robotics}, 2nd ed.,
Monographs in Computer Science, Springer, New York (2005).

\bibitem{Weinberg1995}
S.~Weinberg,
\textit{The Quantum Theory of Fields, Volume I: Foundations},
Cambridge University Press, Cambridge (1995).

\bibitem{Tung1985}
W.-K.~Tung,
\textit{Group Theory in Physics},
World Scientific, Singapore (1985).

\bibitem{FultonHarris1991}
W.~Fulton and J.~Harris,
\textit{Representation Theory: A First Course},
Graduate Texts in Mathematics \textbf{129},
Springer, New York (1991).

\bibitem{Humphreys1972}
J.~E.~Humphreys,
\textit{Introduction to Lie Algebras and Representation Theory},
Graduate Texts in Mathematics \textbf{9},
Springer, New York (1972).

\bibitem{maluf2013}
J.~W.~Maluf,
\textit{The teleparallel equivalent of general relativity},
Annalen der Physik \textbf{525} (5--6), 339--357 (2013).

\end{thebibliography}

\end{document}